% !PRINTED 429
\weberpage{429}

Wir wollen beweisen, dass für jede Primzahl $n$ primitive Wurzeln existiren
und ihre Anzahl bestimmen, und gehen dabei denselben Weg, wie bei den
Einheitswurzeln.

Es seien $a$ und $b$ zwei Theiler von $n-1$, die unter sich relativ prim
sind, und $\alpha$ eine $a^{\mathrm{te}}$, $\beta$ eine $b^{\mathrm{te}}$
primitive Congruenzwurzel von $n$. Setzen wir $ab=c$, so ist
$\gamma=\alpha\beta$ eine primitive $c^{\mathrm{te}}$ Congruenzwurzel, und
umgekehrt lässt sich auch jede primitive $c^{\mathrm{te}}$ Congruenzwurzel
in der Form $\alpha\beta$ darstellen.

Denn es ist erstens: $(\alpha\beta)^c = \alpha^c\beta^c \equiv 1 \pmod{n}$,
und zweitens: wenn $(\alpha\beta)^h = \alpha^h\beta^h \equiv 1 \pmod{n}$
ist, so ist auch $\alpha^{hb}\beta^{hb}\equiv 1$, also $\alpha^{hb}\equiv 1$,
also $hb$ durch $a$ theilbar, also auch $h$ durch $a$ theilbar, und ebenso
schliesst man, dass $h$ durch $b$, also auch durch $ab=c$ theilbar sein
muss; d.\,h. $\alpha\beta$ ist primitive $c^{\mathrm{te}}$
Congruenzwurzel.

Ist umgekehrt $\gamma$ eine primitive $c^{\mathrm{te}}$ Congruenzwurzel, so
bestimmen wir die positiven ganzen Zahlen $x, y$ so, dass
\weq{10}{a y + b x \equiv 1 \pmod{c}}
wird, indem wir zunächst $x$ aus $bx \equiv 1 \pmod{a}$ und dann $y$ aus
$y \equiv \dfrac{1-bx}{a} \pmod{b}$ bestimmen. Dann ist $y$ relativ prim
zu $b$ und $x$ relativ prim zu $a$. Hiernach ist
\[
\gamma \equiv \gamma^{bx}\gamma^{ay} = \alpha\beta \pmod{n}
\]
und $\gamma^{bx} = \alpha$ ist primitive $a^{\mathrm{te}}$
Congruenzwurzel, $\gamma^{ay} = \beta$ primitive $b^{\mathrm{te}}$
Congruenzwurzel (weil $\alpha^h = \gamma^{bxh}$ nur dann $\equiv 1$ sein
kann, wenn $h$ durch $a$ theilbar ist).

Es ist noch zu zeigen, dass die Producte $\alpha\beta$ alle von einander
verschieden sind, d.\,h. dass die Congruenz
\weq{11}{\alpha\beta \equiv \alpha'\beta' \pmod{n},}
wenn $\alpha, \alpha'$; $\beta, \beta'$ primitive $a^{\mathrm{te}}$ und
$b^{\mathrm{te}}$ Congruenzwurzeln sind, nur befriedigt werden kann, wenn
\[
\alpha \equiv \alpha', \quad \beta \equiv \beta' \pmod{n}
\]
ist. Dies folgt aber nach (10), wenn man (11) in die Potenz
\[
bx \equiv 1 \pmod{a}
\]
erhebt, woraus sich $\alpha \equiv \alpha'$ ergiebt, und dann auch
$\beta \equiv \beta'$.

Bezeichnen wir also die Anzahl der primitiven $a^{\mathrm{ten}}$
Congruenzwurzeln mit $\varphi(a)$, so folgt aus dem Bewiesenen
\weq{12}{\varphi(c) = \varphi(a)\varphi(b).}
